\documentclass[10pt,a4paper]{amsart}
\usepackage[margin=19mm]{geometry}
\usepackage{amsmath,amssymb,mathtools}
\usepackage[scr=rsfs,cal=euler]{mathalfa}
\usepackage{xfrac}
\usepackage[unicode,hidelinks,bookmarks=false]{hyperref}
\newcommand{\ie}{i.e.\ }
\newcommand{\cf}{cf.\ }
\newcommand{\resp}{respectively\ }
\newcommand{\Subsection}{\subsection}
\providecommand{\nobreakdash}{\nobreak\hbox{-}}
\setlength{\emergencystretch}{2em}
\multlinegap0pt
\usepackage{amssymb}
\usepackage{float}
\usepackage{bm}
\usepackage{tikz}
\newtheorem{lemma}{Lemma}
\newcommand{\qbinom}[2]{\genfrac{[}{]}{0pt}{}{#1}{#2}}
\title{The DNA Coverage Depth Problem: Duality, \\ Weight Distributions, and Applications}
\author{Matteo Bertuzzo, Alberto Ravagnani, Eitan Yaakobi}
\date{Source: arXiv:2603.06489v1 (CC BY 4.0)}
\begin{document}
\maketitle

\noindent\emph{Source and licence.} The DNA Coverage Depth Problem: Duality, \\ Weight Distributions, and Applications. Version arXiv:2603.06489v1 (CC BY 4.0), distributed by arXiv under the Creative Commons Attribution 4.0 International licence, http://creativecommons.org/licenses/by/4.0/, as recorded in the arXiv licence metadata for this exact version and shown on the abstract page https://arxiv.org/abs/2603.06489v1. Full licence notice: \texttt{licenses/arxiv-2603.06489-CC-BY-4.0.txt}.\par\smallskip

\noindent\textit{Editorial scope.} Counted unit: the article's lemma first\_inversion (source0.tex lines 576--586). The article's complete original proof of it is delivered with it (source0.tex lines 587--604, one \texttt{proof} environment of 18 lines). No mathematics is written, repaired or re-derived: the statement and the proof are the article's own lines, in the article's own notation, and no retained line is substituted or re-flowed.  No further source-local context is needed: every cross-reference of every retained line resolves inside the two delivered spans.  Nothing was omitted from any retained span, so the package carries no omission record.  No retained line contains a citation, so the wrapper carries no bibliography and no bibliography entry.  No retained line contains a figure environment, an image-inclusion command or drawing code, so no image is delivered and the package declares no asset dependency.  Editorial formatting only, in addition: an AMS \texttt{amsart} preamble that restates the article's own declarations and macros used by the retained lines, namely the environments \texttt{lemma} on separate counters, together with the standard packages those lines need.  The retained environments print in this document as Lemma 1 in order of appearance, whereas the article prints the same items with the numbering of its own full document.  The complete native source of the article as distributed in the arXiv e-print for this exact version is bundled unchanged as \texttt{sources/195880/source0.tex}.
\begin{lemma} \label{first_inversion}
Let $\{x_i\}_{i=0}^\infty$ be an integer sequence.
    Let
    \[
    y_m = \sum_{i=0}^m \qbinom{m}{i}_q x_i \quad \textnormal{for all $m \ge 0$}.
    \]
    Then
    \[
    x_i = \sum_{m=0}^i (-1)^{i-m} q^{\binom{i-m}{2}} \qbinom{i}{m}_q y_m \quad \textnormal{ for all $i \ge 0$}.
    \]
\end{lemma}

\begin{proof}
    Fix $i \ge 0$. We have
    \allowdisplaybreaks
    \[
    \begin{split}
    \sum_{m=0}^i (-1)^{i-m} q^{\binom{i-m}{2}} \qbinom{i}{m}_q y_m &= \sum_{m=0}^i (-1)^{i-m} q^{\binom{i-m}{2}} \qbinom{i}{m}_q \sum_{j=0}^m \qbinom{m}{j}_q x_j \\
    &= \sum_{j=0}^i x_j \sum_{m=j}^i (-1)^{i-m} q^{\binom{i-m}{2}} \qbinom{i}{m}_q  \qbinom{m}{j}_q \\
    &= \sum_{j=0}^i x_j \sum_{m=j}^i (-1)^{i-m} q^{\binom{i-m}{2}} \qbinom{i}{j}_q \qbinom{i-j}{i-m}_q \\
    &= \sum_{j=0}^i \qbinom{i}{j}_q x_j \sum_{m=j}^i (-1)^{i-m} q^{\binom{i-m}{2}}  \qbinom{i-j}{i-m}_q \\
    &= \sum_{j=0}^i \qbinom{i}{j}_q x_j \sum_{m=0}^{i-j} (-1)^m q^{\binom{m}{2}}  \qbinom{i-j}{m}_q =  x_i.
    \end{split}
    \]
    The latter identity holds because
    \[
    \sum_{m=0}^{i-j} (-1)^m q^{\binom{m}{2}}  \qbinom{i-j}{m}_q = 0
    \]
    unless $i = j$ (in which case the value is 1), by the $q$-Binomial Theorem.
\end{proof}

\end{document}
